\section{Literature Review}
\label{sec:literature}

The design of an inventory-aware quoting policy originates with
\citet{hostoll1981optimal}, who frame the dealer's problem as a trade-off
between expected spread capture and the variance of an accumulating
inventory. \citet{avellaneda2008high} recast this intuition as a stochastic
control problem with a Poisson fill intensity of the form
$\lambda(\delta) = A\,e^{-k\delta}$ and derive closed-form reservation
prices that tilt quotes against the current position. Because the original
Avellaneda--Stoikov solution depends on a terminal liquidation date that
does not exist for a 24/5 FX book, we follow the infinite-horizon limit
described by \citet{gueant2013dealing} and \citet{gueant2017optimal}, in
which a discount rate $\omega$ replaces the time-to-horizon and the
inventory penalty becomes stationary. This is the formulation implemented
in our Phase~1 quoter.

The second strand of the literature informs our treatment of order flow
and adverse selection. \citet{roll1984simple} gives the effective-spread
decomposition we rely on for cross-venue sanity checks, while
\citet{biais1995empirical} provide the empirical baseline for limit order
book dynamics. \citet{engle2004impacts} and \citet{cont2014price} establish
that short-horizon quote revisions are driven primarily by signed
order-flow imbalance; this motivates our OFI-based skew, which adjusts the
reservation price in the direction of recent imbalance rather than
inferring it from price moves alone. To compensate quoted spreads for the
$150$\,ms latency gap against HFTs, we draw on the adverse-selection and
latency-aware quoting refinements consolidated by
\citet{carteajaimungal2015algorithmic}, which justify the asymmetric
spread regime and the emergency inventory multiplier used in our hedge
controller.

The Phase~2 calibration layer re-estimates the statistical primitives that
Phase~1 took as given. Volatility enters the quoter through the A--S
inventory term $\gamma\sigma^{2}/(2\omega)$; we estimate $\sigma$ with a
GARCH(1,1) model \citep{bollerslev1986generalized} on the simulated
reference mid, and cross-check it against the exponentially weighted
estimator of \citet{riskmetrics1996}, which is better suited to the short
effective horizons over which the quoter reacts. Client order sizes are
heavy-tailed but bounded from above, so we fit a truncated Pareto
distribution following \citet{aban2006parameter}, whose MLE formulation
handles the upper cap cleanly; the power-law regularity of large trades
documented by \citet{gabaix2003theory} is what justifies this parametric
family in the first place.

Relative to this body of work, our contribution is a fully simulated
multi-venue FX designated-market-maker study rather than a new theoretical
model. We combine an infinite-horizon Avellaneda--Stoikov skeleton with an
OFI-based skew, a cross-venue latency-aware fair mid that blends two
reference venues at their respective delays, session-adaptive arrival
decay, and a depth-gated hedge router that externalises inventory on the
reference venues. The whole pipeline is run in two stages: a
heuristic-parameter Phase~1, followed by a Phase~2 in which arrival
intensity, size distribution, volatility, and the risk-aversion $\gamma$
are re-estimated by maximum likelihood on a simulated month of fills.
This two-phase calibration workflow lets us quantify the marginal value
of data-driven parameters over first-principles defaults in an
environment that would otherwise have no trading history at launch.
